% Abstrak berbahasa Indonesia.
\begin{abstrak}
\panduan{Abstrak ditulis satu paragraf utuh (200--250 kata, tanpa sitasi). Memuat latar belakang masalah empiris, tujuan, metode/instrumen, temuan angka utama, dan simpulan. Kata kunci minimal 5 kata/frasa dipisahkan titik koma. Nonaktifkan kotak ini dengan opsi panduan=false.}

Evaluasi kegunaan sistem informasi pelayanan publik penting dilakukan untuk menjamin efektivitas adopsi teknologi serta kepuasan warga di tingkat pemerintahan desa. Penerapan digitalisasi administrasi desa kerap menghadapi kendala berupa keragaman literasi digital masyarakat dan antarmuka sistem yang belum optimal. Penelitian ini bertujuan mengevaluasi tingkat kegunaan aplikasi layanan administrasi kependudukan desa berbasis web menggunakan metode \textit{System Usability Scale} (SUS) serta merumuskan rekomendasi perbaikan antarmuka berbasis bukti empiris. Penelitian menggunakan pendekatan kuantitatif dengan melibatkan 30 orang responden warga desa yang dipilih melalui teknik \textit{purposive sampling}. Pengumpulan data dilakukan menggunakan kuesioner baku SUS 10 butir pertanyaan skala Likert 1--5 yang telah teruji validitas dan reliabilitasnya, dipadukan dengan observasi tugas terstruktur dan wawancara semiterstruktur. Hasil evaluasi menghasilkan nilai rata-rata skor SUS sebesar 74,50 dengan simpangan baku 8,25. Berdasarkan tolok ukur standar SUS, skor ini menempatkan aplikasi pada kategori penerimaan \textit{Acceptable}, predikat kegunaan \textit{Good} pada skala pemeringkatan adjektiva, dan peringkat persentil B, yang mengonfirmasi bahwa sistem dinilai layak secara operasional. Analisis butir instrumen menunjukkan aspek kemudahan integrasi fungsi memperoleh skor tertinggi, sedangkan alur pengunggahan berkas persyaratan dan kebutuhan panduan teknis awal memperoleh penilaian terendah. Berdasarkan temuan tersebut, rekomendasi perbaikan difokuskan pada penyederhanaan tata letak formulir digital, penambahan petunjuk interaktif kontekstual, dan penguatan validasi masukan data secara langsung. Simpulan penelitian menegaskan bahwa aplikasi dinilai efektif dan diterima baik oleh pengguna, dengan potensi reduksi durasi layanan administrasi desa hingga 30\%.

\katakunci{evaluasi kegunaan; sistem informasi; \textit{system usability scale}; administrasi desa; kepuasan pengguna}
\end{abstrak}
